\documentclass[french, 11pt]{article}

\input{setup.tex}

\def\hidetitle{}
\def\pagetitle{Exercice Proba - DS}

\begin{document}

\newif\colles

\begin{center}
    \Large{\bf\pagetitle}
    \vspace*{0.5cm}
\end{center}

\renewcommand*{\E}{\mathbb{E}}
\renewcommand*{\D}{\cursive{D}}

\input{title.tex}

\section*{Exercice.}

Soient $X$ et $Y$ deux variables aléatoires indépendantes à valeurs dans $\N$ définies sur un même espace $(\O,\A,\Pr)$.\\
Pour $|t|<1$, on définit les fonctions génératrices de $X$ et de $Y$ respectivement par :
\begin{equation*}
    G_X(t) = \frac{1}{2-t} \quad \et \quad G_Y(t) = 2 - \sqrt{2-t}
\end{equation*}

\begin{question}{}{}
    Déterminer le développement en série entière de la fonction $G_X$.
    \tcblower
    \vspace*{-0.5cm}
    \begin{equation*}
        G_X(t) = \frac{1}{2-t} = \frac{1}{2}\.\frac{1}{1-\frac{t}{2}} = \frac{1}{2}\sum_{n=0}^{+\infty}\left( \frac{t}{2} \right)^n = \sum_{n=0}^\infty\frac{t^n}{2^{n+1}}.
    \end{equation*}
\end{question}

\begin{question}{}{}
    Donner le terme d'ordre $n\in\N^*$ du développement en série entière de la fonction $t\mapsto(a+t)^{1/2}$.
    \tcblower
    \vspace*{-0.5cm}
    \begin{equation*}
        (1+x)^\a = 1 + \sum_{n=1}^{+\infty}\prod_{i=0}^{n-1}\left( \frac{1}{2} - i \right)\frac{t^n}{n!}
    \end{equation*}
\end{question}

\begin{question}{}{}
    En déduire le développement en série entière de la fonction $G_Y$.
    \tcblower
    \vspace*{-0.5cm}
    \begin{equation*}
        G_Y(t) = 2 - \sqrt{2-t} = 2-\sqrt{2}\sqrt{1-\frac{t}{2}} = 2 - \sqrt{2}\left[ 1 + \sum_{n=1}^{\infty}\binom{1/2}{n}\left( \frac{-t}{2} \right)^n \right] = 2 - \sqrt{2} + \sum_{n=1}^{\infty}\frac{(-1)^{n+1}\sqrt{2}}{2^n}\binom{1/2}{n}t^n
    \end{equation*}
\end{question}

\begin{question}{}{}
    Pour tout $n\in\N$, calculer $\Pr(X=n)$ et $\Pr(Y=n)$.
    \tcblower
    Par définition, $G_X(t) = \sum_{n=0}^\infty t^n \Pr(X=n)$, donc par unicité du DSE, $\Pr(X=n)=\frac{1}{2^{n+1}}$.\\
    De même, $\forall n \in \N, ~ \Pr(Y=0)=2-\sqrt{2}$ et $\Pr(Y=n)=-\frac{\sqrt{2}}{2^nn!}\prod_{i=0}^{n-1}\left( \frac{1}{2}-i \right)$.
\end{question}

\begin{question}{}{}
    Soit $S=X+Y$. Déterminer la loi de $S$.
    \tcblower
    On a $G_S = G_{X+Y} = G_XG_Y$, donc : 
    \begin{align*}
        G_S(t) &= \frac{1}{2-t}(2-\sqrt{2-t}) = \frac{2}{2-t} - \frac{1}{\sqrt{2-t}} = - \frac{1}{\sqrt{2}} + \sum_{n=0}^{\infty}\left( \frac{t}{2} \right)^n - \frac{1}{\sqrt{2}} \sum_{n=1}^{\infty}\binom{-1/2}{n}\frac{(-1)^n}{2^n}t^n\\
        &=1 - \frac{1}{\sqrt{2}} + \sum_{n=1}^{\infty}\frac{1}{2^n}\left[ 1 + \frac{(-1)^{n+1}}{\sqrt{2}} \binom{-1/2}{n} \right]t^n
    \end{align*}
    Par unicité du DSE, $\Pr(S=0) = 1 - \frac{1}{\sqrt{2}}$ et $\Pr(S=n)=\frac{1}{2^n}\left[ 1 + \frac{(-1)^{n+1}}{\sqrt{2}} \binom{-1/2}{n}\right]$
\end{question}

\begin{question}{}{}
    Justifier que la variable aléatoire $X+1$ suit une loi géométrique dont on déterminera le paramètre.
    \tcblower
    On a $\forall n \in \N^*, ~ (X+1=n) = (X=n-1)$ puis $\Pr(X+1=n)=\Pr(X=n-1)=\frac{1}{2^n}=\frac{1}{2}(1-\frac{1}{2})^{n-1}$.\\
    Donc $X+1\sim \cursive{G}(1/2)$.
\end{question}

\begin{question}{}{}
    En déduire l'espérance et la variance de la variable aléatoire $X$.
    \tcblower
    On rappelle que si $Z\sim\cursive{G}(p)$, alors $\E(Z)=\frac{1}{p}$ et $\V(Z)=\frac{1-p}{p^2}$, alors $\E(X+1)=2$ et $\V(X+1)=2$.\\
    $\hra$ $\E(X) = \E(X+1-1) = \E(X+1) - 1 = 1$ par linéarité de l'espérance.\\
    $\hra$ $\V(X) = \V(X+1-1) = \V(X+1) = 2$ car $\forall (a,b) \in \R^2, ~ \V(aX+b) = a^2\V(X)$.
\end{question}

\begin{question}{}{}
    Déterminer à l'aide de la fonction génératrice $G_Y$ l'espérance des variables aléatoires $Y$ et $Y(Y-1)$.
    \tcblower
    On rappelle que si $G_X(t)$ a un rayon de convergence $R>1$, alors $G_X$ est $\cursive{C}^\infty$ sur $]-R,R[$, on peut donc dériver terme-à-terme avec conservation du rayon et par théorème du transfert :
    \begin{equation*}
        G'_X(1) = \E(X) \quad \et \quad G''_X(1) = \E(X(X-1))
    \end{equation*} 
    Donc $\E(Y)=G'_Y(1) = \frac{1}{2}$ et $\E(Y(Y-1))=G''_Y(1)=\frac{1}{4}$.
\end{question}

\pagebreak

\begin{question}{}{}
    En déduire la variance de la variable aléatoire $Y$.
    \tcblower
    Par Koenig-Huygens : $\V(Y) = \E(Y^2) - \E(Y)^2 = \E(Y(Y-1))+\E(Y) - \E(Y)^2 = \frac{1}{4} + \frac{1}{2} - \frac{1}{4} = \frac{1}{2}$. 
\end{question}

\vspace*{-0.5cm}

\begin{question}{}{}
    Déterminer l'espérance et la variance de la variable aléatoire $S$.
    \tcblower
    $\hra$ $\E(S)=\E(X+Y)=\E(X) + \E(Y) = \frac{3}{2}$.\\
    $\hra$ $\V(S)=\V(X+Y)=\V(X) + \V(Y) = \frac{5}{2}$ puisque $X\ind Y$.
\end{question}

\fin

\end{document}
